%   terminology:
%     process: refers to execution in user space, or maybe struct proc &c
%     process memory: the lower part of the address space
%     process has one thread with two stacks (one for in kernel mode and one for
%     in user mode)
% talk a little about initial page table conditions:
%     paging not on, but virtual mostly mapped direct to physical,
%     which is what things look like when we turn paging on as well
%     since paging is turned on after we create first process.
%   mention why still have SEG_UCODE/SEG_UDATA?
%   do we ever really say what the low two bits of %cs do?
%     in particular their interaction with PTE_U
%   sidebar about why it is extern char[]

\chapter{操作系统组织结构}
\label{CH:FIRST}

操作系统的关键需求是同时支持多个活动。例如，可以使用
Chapter~\ref{CH:UNIX}
中的 {\tt fork} 和 {\tt exec} 系统调用（system call）来启动编译器和文本编辑器作为进程（process）。操作系统必须在多个进程之间
\indextext{time-share}
分时（time-sharing）共享 CPU 和内存等资源。操作系统还必须安排
\indextext{isolation}
隔离（isolation）机制，使各进程相互独立。如果一个进程出现错误并发生故障，不应影响其他无关进程。然而，完全隔离又过于严格，因为进程应该能够有意地进行交互；管道（pipeline）就是一个例子。因此，操作系统必须满足三个要求：多路复用（multiplexing）、隔离（isolation）和交互（interaction）。

本章概述了操作系统如何组织以实现这三个要求。实现方式有很多，但本文聚焦于围绕
\indextext{monolithic kernel}
宏内核（monolithic kernel）的主流设计，许多 Unix 操作系统都采用这种架构。本章还概述了 xv6 进程，即 xv6 中的隔离单元。

Xv6 运行在
\indextext{multi-core}
多核（multi-core）RISC-V 微处理器上，
其许多底层功能（例如进程实现）是特定于 RISC-V 的。RISC-V 是 64 位 CPU，xv6 使用 ``LP64'' C 语言编写，这意味着 C 语言中的 long（L）和指针（P）是 64 位，但 {\tt int} 是 32 位。本书假设读者在某些架构上有一定的机器级编程经验，并将在需要时介绍 RISC-V 特定的概念。
用户级 ISA~\cite{riscv:user} 和特权架构~\cite{riscv:priv} 文档是完整的规范。
读者还可以参考
``The RISC-V Reader: An Open Architecture
Atlas''~\cite{riscv}。

完整计算机中的 CPU 被支持硬件包围，其中大部分以 I/O 接口的形式存在。Xv6 为 qemu 的 ``-machine virt'' 选项模拟的支持硬件而编写。这包括 RAM、包含启动代码的 ROM、与用户键盘/屏幕的串行连接，以及用于存储的磁盘。

%%
\section{抽象物理资源}
%%

人们遇到操作系统时可能会问的第一个问题是：为什么需要操作系统？也就是说，可以将
Figure~\ref{fig:api}
中的系统调用实现为一个库，应用程序与之链接。在这种方案中，每个应用程序甚至可以有自己的库，根据其需求定制。应用程序可以直接与硬件资源交互，并以最适合应用的方式使用这些资源（例如，实现高性能或可预测的性能）。一些用于嵌入式设备或实时系统的操作系统就是以这种方式组织的。

这种库方法的缺点是，如果有多个应用程序运行，这些应用程序必须表现良好。例如，每个应用程序必须定期放弃 CPU，以便其他应用程序可以运行。如果所有应用程序相互信任且没有错误，这种
\textit{协作式}
分时方案可能是可行的。更典型的情况是应用程序互不信任，且存在错误，因此通常需要比协作方案更强的隔离。

为了实现强隔离，最好禁止应用程序直接访问敏感的硬件资源，而是将资源抽象为服务。例如，Unix 应用程序仅通过文件系统的
\lstinline{open}、
\lstinline{read}、
\lstinline{write}
和
\lstinline{close}
系统调用与存储交互，而不是直接读写磁盘。这为应用程序提供了路径名的便利，并允许操作系统（提供该接口）管理磁盘。即使隔离不是关注点，有意交互（或只是希望互不干扰）的程序也可能发现文件系统是比直接使用磁盘更方便的抽象。

同样，Unix 透明地在进程之间切换硬件 CPU，根据需要保存和恢复寄存器状态，使应用程序无需感知分时。这种透明性允许操作系统共享 CPU，即使某些应用程序处于无限循环中。

再举一个例子，Unix 进程使用
\lstinline{exec}
构建其内存映像，而不是直接与物理内存交互。这允许操作系统决定将进程放置在内存的何处；如果内存紧张，操作系统甚至可能将进程的某些数据存储在磁盘上。
\lstinline{exec}
还为用户提供了使用文件系统存储可执行程序映像的便利。

Unix 进程之间的许多交互形式通过文件描述符（file descriptor）进行。文件描述符不仅抽象了许多细节（例如，管道或文件中的数据存储位置），而且其定义方式也简化了交互。例如，如果管道中的一个应用程序退出或失败，内核会自动为管道中的下一个进程生成文件结束信号。

Figure~\ref{fig:api}
中的系统调用接口经过精心设计，既提供了程序员便利，又提供了强隔离的可能性。Unix 接口不是抽象资源的唯一方式，但已被证明是一种良好的方式。

%%
\section{用户模式、监管者模式和系统调用}
%%

强隔离需要在应用程序和操作系统之间建立严格的边界。即使应用程序有错误或恶意行为，也不应允许其干扰操作系统或其他程序的运行。为了实现强隔离，操作系统必须确保应用程序不能修改（甚至不能读取）操作系统的数据结构和指令，并且应用程序不能访问其他进程的内存。

CPU 为强隔离提供硬件支持。例如，RISC-V 有三个特权级别（privilege level），限制代码可以执行的操作：
\indextext{machine mode}（机器模式）、
\indextext{supervisor mode}（监管者模式）和
\indextext{user mode}（用户模式）。在机器模式下执行的指令拥有完全特权；CPU 启动时处于机器模式。机器模式主要用于在启动期间配置计算机。Xv6 在机器模式下执行一小段时间，然后切换到监管者模式。

在监管者模式下，CPU 被允许执行
\indextext{privileged instructions}（特权指令）：例如，启用和禁用中断、读写保存页表地址的寄存器等。如果用户模式中的应用程序尝试执行特权指令，CPU 不会执行该指令，而是会``陷入''（trap）到监管者模式中的特殊代码，该代码可以终止应用程序。
Chapter~\ref{CH:UNIX}
中的 Figure~\ref{fig:os} 说明了这种组织结构。应用程序只能执行用户模式指令（例如，数字相加等），被称为在
\indextext{user space}（用户空间）
中运行，而监管者模式中的软件还可以执行特权指令，被称为在
\indextext{kernel space}（内核空间）
中运行。在内核空间（或监管者模式）中运行的软件称为
\indextext{kernel}（内核）。

应用程序通过系统调用（system call）与内核交互，例如 {\tt read}。应用程序不允许直接调用内核函数或访问内核内存。RISC-V 提供了 \indexcode{ecall} 指令用于系统调用；它将 CPU 从用户模式切换到监管者模式，并跳转到内核指定的入口点。一旦 CPU 切换到监管者模式，内核就可以验证系统调用的参数（例如，检查传递给系统调用的地址是否是应用程序内存的一部分），决定应用程序是否被允许执行请求的操作（例如，检查应用程序是否被允许写入指定文件），然后拒绝或执行该操作。内核控制进入监管者模式的入口点非常重要；如果应用程序可以决定内核入口点，恶意应用程序可能会在内核跳过参数验证的位置进入，例如。

%%
\section{内核组织结构}
%% 

一个关键的设计问题是操作系统的哪部分应该在监管者模式下运行。一种可能是整个操作系统驻留在内核中，这样所有系统调用的实现都在监管者模式下运行。这种组织称为
\indextext{monolithic kernel}（宏内核）。

在宏内核组织中，整个操作系统由一个在监管者模式下运行的单一程序组成。这种组织方便的一个原因是操作系统设计者不需要将代码划分为需要监管者特权和不需要的部分。此外，操作系统不同部分之间的协作很容易，因为它们是同一程序的一部分。例如，宏内核可能与文件系统和虚拟内存系统共享磁盘块缓存。

缺点是宏内核往往变得庞大而复杂，以至于没有一个开发者理解代码不同部分之间的所有交互；这是产生错误的根源。内核中的错误特别麻烦，因为它可能导致整个计算机崩溃，或导致许多应用程序故障，或使整个计算机容易受到安全攻击。

\indextext{microkernel}（微内核）
旨在减少内核中错误的发生。其想法是将绝对最少的功能放入内核本身，这样少量代码在监管者模式下执行，内核易于理解和分析其正确性。操作系统的大部分作为用户级服务器进程运行。例如，文件系统代码将作为服务器进程执行，在用户模式而非监管者模式下。

\begin{figure}[t]
\center
\includegraphics[scale=0.5]{fig/mkernel.pdf}
\caption{A microkernel with a file-system server}
\label{fig:mkernel}
\end{figure}

Figure~\ref{fig:mkernel}
说明了这种微内核设计。在图中，文件系统作为用户级服务器进程运行。为了允许应用程序与文件服务器交互，内核提供了一种进程间通信（inter-process communication）机制，用于将消息从一个用户模式进程发送到另一个用户模式进程。例如，如果像 shell 这样的应用程序想要读取或写入文件，它会向文件服务器发送消息并等待响应。

在微内核中，内核接口由一些低级函数组成，用于启动应用程序、发送消息、访问设备硬件等。这种组织允许内核相对简单，因为大部分操作系统驻留在用户级服务器中。

在现实世界中，宏内核和微内核都很流行。许多 Unix 内核是宏内核。例如，Linux 具有宏内核，尽管某些操作系统功能作为用户级服务器运行（例如，窗口系统）。Linux 为操作系统密集型应用程序提供高性能，部分原因是内核的子系统可以紧密集成。

Minix、L4 和 QNX 等操作系统被组织为带有服务器的微内核，并已在嵌入式环境中广泛部署。L4 的一个变体 seL4 足够小，已被验证具有内存安全性和其他安全属性~\cite{sel4}。

关于哪种组织更好的争论在操作系统开发者之间存在很大争议，但没有确凿证据表明哪一种更好。此外，这很大程度上取决于``更好''的含义：更快的性能、更小的代码大小、内核的可靠性、完整操作系统的可靠性（包括用户级服务）等。

还有一些实际考虑可能比组织问题更重要。某些操作系统具有微内核，但出于性能原因，某些用户级服务在内核空间中运行。某些操作系统具有宏内核，因为它们就是这样开始的，并且几乎没有动力转向纯微内核组织，因为新功能可能比重写现有操作系统以适应微内核设计更重要。

从本书的角度来看，微内核和宏内核操作系统共享许多关键思想。它们实现系统调用，使用页表（page table），处理中断（interrupt），支持进程，使用锁（lock）进行并发控制，实现文件系统（file system）等。本书专注于这些核心思想。

Xv6 被实现为宏内核，与大多数 Unix 操作系统一样。因此，xv6 内核接口对应于操作系统接口，内核实现完整的操作系统。由于 xv6 不提供许多服务，其内核比某些微内核更小，但从概念上讲 xv6 是宏内核。  

\section{代码：xv6 组织结构}

\begin{figure}[t]
\center
\begin{tabular}{l|l|l}
& {\bf 文件} & {\bf 描述} \\
\midrule
启动 & entry.S & 最初的启动指令。 \\
 & main.c & 控制其他模块的初始化。 \\
 & start.c & 早期的机器模式启动代码。 \\
进程 & exec.c & exec() 系统调用。 \\
 & proc.c & 进程和调度。 \\
 & swtch.S & 线程切换。 \\
 & sysproc.c & 进程相关的系统调用。 \\
陷入 & kernelvec.S & 处理来自内核代码的陷入。 \\
 & trampoline.S & 处理来自用户代码的陷入。 \\
 & trap.c & 处理和从陷入与中断返回的 C 代码。 \\
 & syscall.c & 将系统调用分派给处理函数。 \\
内存 & vm.c & 管理页表和地址空间。 \\
 & kalloc.c & 物理页分配器。 \\
设备 & console.c & 连接用户键盘和屏幕。 \\
 & plic.c & RISC-V 中断控制器。 \\
 & printf.c & 格式化输出到控制台。 \\
 & uart.c & 串口控制台设备驱动。 \\
 & virtio\_disk.c & 磁盘设备驱动。 \\
文件系统 & bio.c & 文件系统的磁盘块缓存。 \\
 & file.c & 文件描述符支持。 \\
 & fs.c & 文件系统。 \\
 & log.c & 文件系统日志和崩溃恢复。 \\
 & sysfile.c & 文件相关的系统调用。 \\
 & pipe.c & 管道。 \\
其他 & sleeplock.c & 让出 CPU 的锁。 \\
 & spinlock.c & 不让出 CPU 的锁。 \\
 & string.c & C 字符串和字节数组库。 \\
\end{tabular}
\caption{Xv6 内核源文件。}
\label{fig:source}
\end{figure}

xv6 内核源代码位于 {\tt kernel/} 子目录中。
Figure~\ref{fig:source} 列出了这些文件，按内核的主要职责领域划分：启动系统（booting）、创建和控制进程、处理陷入（中断和系统调用）、分配内存和配置虚拟地址、控制设备，以及管理文件系统。

%% 
\section{Process overview}
%% 

The unit of isolation in xv6 (as in other Unix operating systems) is a 
\indextext{process}.
The process abstraction prevents one process from wrecking or spying on
another process's memory, CPU, file descriptors, etc.  It also prevents a process
from wrecking the kernel itself, so that a process can't subvert the kernel's
isolation mechanisms.
The kernel must implement the process abstraction with care because
a buggy or malicious application may trick the kernel or hardware into doing
something bad (e.g., circumventing isolation).  The mechanisms used by
the kernel to implement processes include the user/supervisor mode flag, address spaces,
and time-slicing of threads.

To help enforce isolation, the process abstraction provides the
illusion to a program that it has its own private machine.  A process provides
a program with what appears to be a private memory system, or
\indextext{address space}, 
which other processes cannot read or write.
A process also provides the program with what appears to be its own
CPU to execute the program's instructions.

Xv6 uses page tables (which are implemented by hardware) to give each process
its own address space. The RISC-V page table
translates (or ``maps'') a
\indextext{virtual address}
(the address that an RISC-V instruction manipulates) to a
\indextext{physical address}
(an address that the CPU sends to main memory).

\begin{figure}[t]
\centering
\includegraphics[scale=0.5]{fig/as.pdf}
\caption{Layout of a process's virtual address space}
\label{fig:as}
\end{figure}

Xv6 maintains a separate page table for each process that defines that process's
address space. As illustrated in 
Figure~\ref{fig:as},
an address space includes the process's
\indextext{user memory}
starting at virtual address zero. Instructions come first,
followed by global variables, then the stack,
and finally a ``heap'' area (for malloc)
that the process can expand as needed.
There are a number of factors that limit the
maximum size of a process's address space:
pointers on the RISC-V are 64 bits wide;
the hardware uses only the low 39 bits when
looking up virtual addresses in page tables;
and xv6 uses only 38 of those 39 bits.
Thus, the maximum address is $2^{38}-1$ =
0x3fffffffff, which is \lstinline{MAXVA}~\lineref{kernel/riscv.h:/define.MAXVA/}.
At the top of the address space xv6 places
a \indextext{trampoline} page (4096 bytes)
and a \indextext{trapframe} page.  Xv6 uses these two pages
to transition into the kernel and back;
the trampoline page contains the code to transition in and out
of the kernel, and the trapframe is where the kernel saves
the process's user registers,
as Chapter~\ref{CH:TRAP} explains.

The xv6 kernel maintains many pieces of state for each process,
which it gathers into a
\indexcode{struct proc}
\lineref{kernel/proc.h:/^struct.proc/}.
A process's most important pieces of kernel state are its 
page table, its kernel stack, and its run state.
We'll use the notation
\indexcode{p->xxx}
to refer to elements of the
\lstinline{proc}
structure; for example,
\indexcode{p->pagetable} is a pointer to the process's page table.

At this point, please read {\tt kernel/proc.h}, which defines {\tt
  struct proc}. The xv6 code is more important for you to understand
than this book; you should prioritize the code, and consult this book
as needed to clarify the code. The purpose of some of the code may not
be apparent at first, but further reading and searching the code will
help. Feel free to explore and modify the code.

Each process has a thread of control (or 
\indextext{thread}
for short) that holds the state needed to execute the process.
At any given time, a thread might be executing on a CPU,
or suspended (not executing, but capable of resuming executing
in the future).
To switch a CPU between processes,
the kernel suspends the thread currently running on that CPU
and saves its state,
and restores the state of another process's previously-suspended
thread.  Much of the state of a thread (local variables, function call return
addresses) is stored on the thread's stacks.
Each process has two stacks: a user stack and a kernel stack
(\indexcode{p->kstack}).
When the process is executing user instructions, only its user stack
is in use, and its kernel stack is empty.
When the process enters the kernel (for a system call or interrupt),
the kernel code executes on the process's kernel stack; while
a process is in the kernel, its user stack still contains saved
data, but isn't actively used.
A process's thread alternates between actively using its user stack
and its kernel stack. The kernel stack is separate (and protected from
user code) so that the kernel
can execute even if a process has wrecked its user stack.

A process's user code
can make a system call by executing the RISC-V \indexcode{ecall}
instruction. This instruction switches to supervisor mode and
changes the program counter to a kernel-defined entry point.  The code
at the entry point switches to the process's kernel stack and executes the kernel
instructions that implement the system call.  When the system call
completes, the kernel returns to
user space by executin the \indexcode{sret} instruction, which switches
to user mode and resumes executing user instructions
just after the system call instruction.  A process's thread can
``block'' in the kernel to wait for I/O, and resume where it left off
when the I/O has finished.

\indexcode{p->state} 
indicates whether the process is allocated, ready
to run, currently running on a CPU, waiting for I/O, or exiting.

\indexcode{p->pagetable}
holds the process's page table, in the format
that the RISC-V hardware expects.
Xv6 causes the paging hardware to use a process's
\lstinline{p->pagetable}
when executing that process in user space.
A process's page table also serves as the record of the
addresses of the physical pages allocated to store the process's
memory.

In summary, a process bundles two design ideas: an address space to
give a process the illusion of its own memory, and a thread to give
the process the illusion of its own CPU.  In xv6, a process consists
of one address space and one thread.  In real operating systems a
process may have more than one thread to take advantage of multiple CPUs.

%% 
\section{Code: starting xv6, the first process and system call}
%% 
To make xv6 more concrete, we'll outline how the kernel starts and
runs the first process. The subsequent chapters will describe the
mechanisms that show up in this overview in more detail.
Please read {\tt kernel/entry.S}, {\tt kernel/start.c},
{\tt kernel/main.c}, and {\tt user/init.c}.

When the RISC-V computer powers on, it initializes
itself and runs a boot loader which is stored in read-only
memory.  The boot loader copies the xv6 kernel into memory
at physical address
\texttt{0x80000000}.
The reason it places the kernel at
\texttt{0x80000000}
rather than
\texttt{0x0}
is because the address range
\texttt{0x0:0x80000000}
contains I/O devices.

Then the boot loader jumps to xv6 starting at
\indexcode{_entry}
\lineref{kernel/entry.S:/^.entry:/}.
The RISC-V starts with paging hardware disabled:
virtual addresses map directly to physical addresses.
The instructions at
\lstinline{_entry}
set up a stack so that xv6 can run C code.
Xv6 declares space for this stack,
\lstinline{stack0},
in the file
\lstinline{start.c}
\lineref{kernel/start.c:/stack0/}.
The code at
\lstinline{_entry}
loads the stack pointer register
\texttt{sp}
with the address
\lstinline{stack0+4096},
the top of the stack, because the stack
on RISC-V grows down.
Now that the kernel has a stack,
\lstinline{_entry}
calls into C code at
\lstinline{start}
\lineref{kernel/start.c:/^start/}.

The function
\lstinline{start}
performs some setup that the CPU only allows in machine mode,
most crucially programming the clock chip to generate
timer interrupts.
Then {\tt start} uses the RISC-V {\tt mret}
instruction to switch to supervisor mode and
jump to 
\lstinline{main}
\lineref{kernel/main.c:/^main/}.
{\tt mret} requires
a bit of setup:
{\tt start} sets the previous privilege mode to
supervisor in the register
\lstinline{mstatus},
sets the destination address to
\lstinline{main}
by writing
\lstinline{main}'s
address into
the register
\lstinline{mepc},
disables virtual address translation in supervisor mode
by writing
\lstinline{0}
into the page-table register
\lstinline{satp},
and delegates all interrupts and exceptions
to supervisor mode.

After
\lstinline{main}
\lineref{kernel/main.c:/^main/}  
initializes several devices and subsystems, 
it creates the first process by calling 
\lstinline{userinit}
\lineref{kernel/proc.c:/^userinit/}.
All newly created processes start executing in the kernel in
\lstinline{forkret}
\lineref{kernel/proc.c:/^forkret/}.
As a special case for the first process, {\tt forkret}
calls {\tt kexec} to load the user program {\tt /init}.

After calling
\lstinline{kexec},
{\tt forkret} returns to user space in
the \lstinline{/init} process.
\lstinline{init}
\lineref{user/init.c:/^main/}
creates a new console device file
if needed
and then opens it as file descriptors 0, 1, and 2.
Then it starts a shell on the console.
The system is up.

\section{Security Model}

You may wonder how the operating system deals with buggy or malicious
code. Because coping with malice is strictly harder than dealing with
accidental bugs, it's reasonable to 
focus mostly on providing security against malice.
Here's a high-level view of typical security assumptions and
goals in operating system design.

The operating system must assume that a process's user-level code will
do its best to wreck the kernel or other processes. User code may try
to dereference pointers outside its allowed address space; it may
attempt to execute instructions not intended
for user code; it may try to read and write RISC-V control
registers; it may try to access device hardware;
and it may pass clever values to system calls in an attempt
to trick the kernel into crashing or doing something stupid.

The
kernel's goal is to restrict each user process so that it can only
access its own user memory, use the 32 general-purpose
RISC-V registers, and affect the kernel and other processes in the
ways that system calls are intended to allow. The kernel must prevent
any other actions. These are typically absolute requirements in kernel
design.

Expectations for the kernel's own code are different. Kernel
code is assumed to be written by well-meaning and careful programmers,
to be bug-free, and to contain
nothing malicious. This assumption affects how we analyze kernel code.
For example, there are many internal kernel functions (e.g., the spin
locks) that would cause serious problems if kernel code used them
incorrectly. We assume,
however, that the kernel uses its own functions correctly.
At the hardware level, the RISC-V CPU, RAM, disk, etc. are
assumed to operate as advertised in the documentation, with no
hardware bugs.

Real life is not so straightforward. It's
difficult to prevent abusive user programs from 
calling system calls in a way that makes the system unusable
by consuming kernel-protected resources:
disk space, CPU time, process table slots, etc. It's usually
impossible to write 100\% bug-free kernel code or design bug-free hardware; if the
writers of malicious user code are aware of kernel or hardware bugs,
they will exploit them. Even in mature, widely-used kernels, such as
Linux, people often discover previously-unknown
vulnerabilities~\cite{mitre:cves}.
Finally, the distinction between
user and kernel code is sometimes blurred: some privileged user-level
processes may provide essential services and effectively be part of
the operating system, and in some operating systems privileged user
code can insert new code into the kernel (as with Linux's loadable
kernel modules and eBPF).

As a partial defense against kernel bugs, xv6 code includes checks for
inconsistencies and unrecoverable errors, and will ``panic'' in
response, by calling {\tt panic()}. This function prints an error
message and halts the system. Panicking is not desirable, but is
preferable to continuing execution. Typically a panic results from a
kernel bug that causes kernel data to be incorrect or causes the
kernel to perform an illegal action such as referencing non-existent
memory; in such a situation it is safer to halt execution with {\tt
  panic()} than to try to continue in an inconsistent state. A kernel
developer would react to a panic by working to identify and fix the
underlying code bug.

%% 
\section{Real world}
%% 

Most operating systems have adopted the process concept, and most
processes look similar to xv6's.  Modern operating systems, however,
support several threads within a process, to allow a single process to
exploit multiple CPUs.  Supporting multiple threads in a
process involves quite a bit of machinery that xv6 doesn't have,
often including interface changes (e.g., Linux's
\lstinline{clone},
a variant of
\lstinline{fork}),
to control which aspects of
a process threads share.
%% 
\section{Exercises}
%%

\begin{enumerate}

\item Add a system call to xv6
  that returns the amount of free memory available.

% break *0x3ffffff000
% disas 0x3ffffff000,+8
% set disassemble-next-line auto
% x/i $pc
% break *0x3ffffff10e
% print/x $sepc
% print/x $pc

\end{enumerate}
